% TOP_PROBLEM: {"id": 7000003, "title": "Geometry of Curves and Surfaces — Problem 1.3", "category_id": 6, "category": {"id": 6, "name": "geometry", "display_name": "Geometry", "description": "Euclidean and non-Euclidean geometry, geometric structures.", "slug": "geometry", "order_index": 6, "created_at": "2026-07-31T15:26:25.671Z"}, "status": "open", "problem_number": "AMR-069-0003", "set_id": 13}
% FIRST_POSTED: 2026-10-01
% ENTRY_KIND: statement_only
% UnsolvedMath ID 7000003; source catalogue code AMR-069-0003
% !TeX program = lualatex
% Definition and sources only; this file is not an LLM solution attempt.
\documentclass[11pt,a4paper]{article}
\usepackage[margin=25mm]{geometry}
\usepackage{fontspec}
\setmainfont{lmroman10-regular.otf}[BoldFont=lmroman10-bold.otf,
 ItalicFont=lmroman10-italic.otf,BoldItalicFont=lmroman10-bolditalic.otf]
\usepackage{amsmath,amssymb,mathtools}
\usepackage{unicode-math}
\setmathfont{latinmodern-math.otf}
\AtBeginDocument{\let\setminus\smallsetminus}
\usepackage{xurl,hyperref}
\hypersetup{colorlinks=true,urlcolor=blue}
\setcounter{secnumdepth}{0}
\setlength{\parindent}{0pt}
\setlength{\parskip}{5pt}
\setlength{\emergencystretch}{3em}
\begin{document}
\section*{Geometry of Curves and Surfaces — Problem 1.3}\label{7000003}
{\small UnsolvedMath ID 7000003 · AMR-069-0003 · Geometry · AMR Open Problem Lists}
\subsection{Definitions and mathematical statement}
Let $A$ be a compact smooth annulus with a Riemannian metric $g$
of strictly negative Gauss curvature in its interior. Prescribe
two disjoint smooth simple closed convex planar curves
$\Gamma_0,\Gamma_1\subset\mathbb R^3$ as its boundary components.
Each curve lies in a plane; the planes need not be the same.
Consider smooth embeddings $f:A\to\mathbb R^3$ inducing the metric
$g$ and carrying the two boundary circles onto these fixed curves.

The rigidity question asks whether two such isometric realizations
with the same boundary data must be congruent by an ambient
Euclidean isometry respecting that data. In a marked formulation,
fix the boundary parametrizations as well and compare embeddings
that agree there. In particular, rigidity would prohibit a
nontrivial continuous isometric deformation with the boundary
held fixed.

The negative curvature concerns the intrinsic metric of the
annulus, equivalently the product of its two principal curvatures
under an isometric embedding. Fixing the boundary means preserving
the actual space curves, not only their lengths or their planar
congruence classes. The problem is about intrinsic isometries of
surfaces; having the same Gauss curvature function alone is a
weaker requirement and is not the proposed rigidity hypothesis.
\subsection{Short English statement}
Is a negatively curved annulus uniquely determined, up to rigid
motion, by its intrinsic metric and its two fixed convex planar
boundary curves?
\subsection{Sources}
\begin{itemize}
\item[S1] ULAM.AI and the UnsolvedMath Contributors, supplied problems.json, record 7000003 (AMR-069-0003), AMR Open Problem Lists. Catalogue identity and source transcription; CC BY 4.0.
\url{https://huggingface.co/datasets/ulamai/UnsolvedMath}.
\item[S2] Ghomi - Open Problems in Geometry of Curves and Surfaces (2019), Problem 1.3, PDF page 6. Original source indexed by AMR-069-0003.
\url{https://people.math.gatech.edu/~ghomi/Papers/op.pdf}.
\end{itemize}
\end{document}
